\chapter{Cervical Infection}
\index{Cervical Infection}

\section*{Definition and Overview}
Cervical infection, medically known as cervicitis, is inflammation of the cervix, the lower part of the uterus that opens into the vagina. It may be caused by infectious organisms, most commonly sexually transmitted infections such as chlamydia and gonorrhoea, or by non-infectious irritation. In many cases it produces no symptoms at all, which is why routine screening is important.

Cervicitis is very common and often resolves with treatment, but untreated infections can spread upward into the uterus and fallopian tubes, causing pelvic inflammatory disease, chronic pelvic pain, and fertility problems. Prompt diagnosis and treatment, together with partner treatment, prevent these complications.

\section*{Epidemiology}
Cervicitis is one of the most common reasons for gynaecological consultation, and sexually transmitted infections (STIs) are among the most frequently diagnosed infections worldwide. Chlamydia, in particular, is extremely common in sexually active young people under the age of 25, and most affected women have no symptoms.

The prevalence of gonorrhoea varies by region and population, with higher rates in urban areas and among certain risk groups. Rates of STIs have been rising in many countries in recent decades, partly because of increased testing and partly because of reduced condom use. Because many cases are silent, the true burden of cervical infection is higher than reported figures suggest.

\section*{Aetiology and Pathophysiology}
Cervicitis develops when organisms or irritants cause inflammation of the cervix. The main causes are:

\begin{itemize}
    \item \textbf{\textit{Chlamydia trachomatis}:} the most common infectious cause, particularly in young women
    \item \textbf{\textit{Neisseria gonorrhoeae}:} a common cause of more severe infection
    \item \textbf{\textit{Mycoplasma genitalium}:} a recognised cause of cervicitis and urethritis
    \item \textbf{Herpes simplex virus:} causes painful, ulcerated cervical and vaginal lesions
    \item \textbf{Trichomonas vaginalis:} a parasitic infection that also causes vaginal discharge and irritation
    \item \textbf{Other bacteria:} such as Group B streptococcus and organisms associated with bacterial vaginosis
    \item \textbf{Non-infectious causes:} irritants including douching, tampon use, chemical contraceptives, and allergic reactions
\end{itemize}

The organisms infect the glandular cells lining the cervix, triggering an inflammatory response with redness, swelling, discharge, and in some cases surface ulceration or bleeding. The infection can extend through the cervical opening into the upper genital tract if left untreated.

\section*{Clinical Features}
Many women with cervicitis have no symptoms. When symptoms occur, they may include:

\begin{itemize}
    \item Abnormal vaginal discharge, which may be watery, cloudy, or purulent
    \item Bleeding after intercourse or between periods
    \item Pain or discomfort during intercourse
    \item Irritation, itching, or burning in the vaginal area
    \item A dull ache or discomfort in the lower abdomen
    \item Pain or burning on passing urine if the urethra is also affected
    \item In herpes infection: painful genital sores, fever, and flu-like symptoms
\end{itemize}

\section*{Differential Diagnosis}
\begin{itemize}
    \item Vaginitis, including thrush (candida), bacterial vaginosis, and trichomoniasis, which affect the vagina
    \item Cervical ectropion or erosion, a normal exposure of glandular tissue that can bleed on contact
    \item Cervical polyps, which are benign growths that may cause bleeding
    \item Early cervical cancer or pre-cancerous change, which can present with discharge and bleeding
    \item Pelvic inflammatory disease, where infection has spread into the uterus and fallopian tubes
    \item Bleeding related to pregnancy, contraception, or the menstrual cycle
\end{itemize}

\section*{When to Seek Medical Care}
See a doctor or sexual health clinic if you have abnormal discharge, bleeding after sex, pain with intercourse, or any other symptoms that concern you. Testing is also recommended after unprotected sex with a new partner, and as part of routine sexual health screening.

\textbf{Contact your local emergency services for:} severe lower abdominal or pelvic pain, fever with chills, heavy bleeding, or symptoms that suggest pelvic inflammatory disease with spreading infection.

\section*{Diagnosis and Investigations}
\begin{itemize}
    \item \textbf{History and examination:} discussion of symptoms, sexual history, contraception, and a speculum examination of the cervix and vagina
    \item \textbf{Swabs:} samples taken from the cervix and vagina and tested by PCR for chlamydia, gonorrhoea, trichomonas, and mycoplasma
    \item \textbf{Urine testing:} a first-void urine sample can also be tested for chlamydia and gonorrhoea
    \item \textbf{Blood tests:} including HIV and syphilis screening where appropriate
    \item \textbf{Colposcopy:} a magnified examination of the cervix if the appearance is unusual or a smear is abnormal
    \item \textbf{Pelvic examination:} checking for tenderness of the uterus and fallopian tubes, which suggests upward spread
\end{itemize}

\section*{Management}
\begin{itemize}
    \item \textbf{Antibiotics:} chlamydia is treated with a course of antibiotics such as doxycycline or azithromycin; gonorrhoea usually needs an injection plus oral treatment
    \item \textbf{Partner treatment:} current and recent sexual partners should be tested and treated to prevent reinfection
    \item \textbf{Abstinence during treatment:} avoiding sexual intercourse until the course is finished and symptoms have settled
    \item \textbf{Antiviral treatment:} for herpes infection, antiviral medicines reduce the duration and severity of episodes
    \item \textbf{Follow-up testing:} a repeat test several weeks after treatment to confirm the infection has cleared
    \item \textbf{Treatment of complications:} pelvic inflammatory disease needs a longer course of antibiotics and close follow-up
    \item \textbf{Screening:} routine cervical screening continues to be important alongside STI testing
\end{itemize}

\section*{Complications}
\begin{itemize}
    \item Pelvic inflammatory disease, with infection spreading to the uterus, fallopian tubes, and ovaries
    \item Scarring and blockage of the fallopian tubes, which can cause infertility
    \item Increased risk of ectopic pregnancy after tubal damage
    \item Chronic pelvic pain
    \item Recurrent or persistent infection if partners are not treated
    \item In pregnancy, untreated infection can be passed to the baby or cause complications
    \item Psychological effects, including anxiety about sexual health and fertility
\end{itemize}

\section*{Prognosis and Outcomes}
With prompt antibiotic treatment, most cervical infections clear completely and cause no lasting harm. Outcomes depend on how quickly treatment starts: infection that has already spread to the upper genital tract carries a higher risk of infertility and chronic pain. Regular screening and partner treatment greatly reduce recurrence. For most women diagnosed early, fertility and long-term health are unaffected.

\section*{Prevention and Self-Care}
\begin{itemize}
    \item Use condoms consistently, which reduce the transmission of most sexually transmitted infections
    \item Have regular sexual health screening, especially if you are under 25 or have a new partner
    \item Limit the number of sexual partners and discuss sexual health with partners
    \item Attend for testing promptly after unprotected sex or if symptoms develop
    \item Complete prescribed antibiotic courses fully, even if symptoms improve
    \item Ensure sexual partners are tested and treated at the same time
    \item Avoid douching, which can disturb normal vaginal flora and increase infection risk
\end{itemize}

\section*{Sources and Further Reading}
The following organisations provide reliable information on cervical infection and sexually transmitted infections.

\begin{itemize}
    \item National Health Service. \textit{NHS conditions guide on chlamydia, gonorrhoea, and sexually transmitted infections.} https://www.nhs.uk/conditions/
    \item Mayo Clinic. \textit{Cervicitis: symptoms, causes, and treatment.} https://www.mayoclinic.org/
    \item Centers for Disease Control and Prevention. \textit{Sexually transmitted infections and their treatment.} https://www.cdc.gov/
    \item World Health Organization. \textit{Global guidance on sexually transmitted infections.} https://www.who.int/
    \item MedlinePlus. \textit{Consumer health information on cervical and vaginal infections.} https://medlineplus.gov/
    \item MSD Manual. \textit{Cervicitis -- professional overview.} https://www.msdmanuals.com/
\end{itemize}
